\begin{abstract}
We present \parbench{}, a kernel-centric benchmark framework for evaluating LLM-based parallel code translation. \parbench{} combines declarative benchmark specifications, an end-to-end build-run-verify harness, survey-grounded corpus construction from 35~open-source HPC repositories, and AST-driven semantics-preserving source augmentation. The corpus presented in this paper contains 96~specifications (87 verified-pass, 9 excluded for platform incompatibilities) over 35~kernels from five benchmark suites, yielding six standard translation directions plus four OpenMP Offload case-study directions.
% src: paper_data_*.json > passk_campaign > aggregate_passk; statistical_analysis.json > model_comparison > pairwise[0]
Under stochastic sampling (temperature~0.7, three independent attempts per task), pass@1 ranges from 23.9\% (\qwenshort{}) to 62.7\% (shared by \gptnew{} and \codexshort{}) over 142~unique source-target pairs, with the code-specialized \codexshort{} performing indistinguishably from the general-purpose \gptnew{} ($p = 1.0$, OR\,=\,0.93). Build-stage adaptation failures account for 50.0\% of \qwenshort{} L0 outcomes, making incomplete API-surface mapping the dominant bottleneck. On L0-conditional augmentation subsets, pass rates show no evidence of memorization-dependent degradation for the GPT models. \parbench{} exposes both current capability and the structural barriers---direction asymmetry, multi-file coordination, and incomplete API adaptation---that remain for reliable parallel code translation.
\end{abstract}